\documentclass[11pt,a4paper]{article}
\usepackage[T1]{fontenc}\usepackage{lmodern}
\usepackage[margin=26mm]{geometry}
\usepackage{amsmath,amssymb,amsthm,mathtools,microtype,booktabs}
\usepackage[hidelinks]{hyperref}\usepackage{xurl}
\hypersetup{pdftitle={Infinite Taylor Sign Changes in Integer Thue Morse Pressure},pdfauthor={Research note prepared for Vladimir Reshetnikov with OpenAI}}
\setlength{\parskip}{3pt}
\newtheorem{theorem}{Theorem}[section]\newtheorem{lemma}[theorem]{Lemma}
\newtheorem{proposition}[theorem]{Proposition}\newtheorem{corollary}[theorem]{Corollary}
\newcommand{\T}{\mathbb T}\newcommand{\C}{\mathbb C}
\newcommand{\R}{\mathbb R}\newcommand{\E}{\mathbb E}
\newcommand{\Lop}{\mathcal L}
\title{Infinite Taylor Sign Changes in\\Integer Thue--Morse Pressure}
\author{Research note prepared for Vladimir Reshetnikov with OpenAI}
\date{1 October 2026}
\begin{document}\maketitle
\begin{abstract}
For every fixed integer $m\ge2$, the Taylor series of integer Thue--Morse pressure at the atomic phase has infinitely many strictly positive and infinitely many strictly negative even coefficients.
Both sign subsequences attain the full exponential coefficient scale.
In fact, each of the degree classes $0$ and $2$ modulo $4$ separately contains infinitely many coefficients of both signs.
The proof combines the repository's real-phase simple Perron theorem with an elementary complex matrix bound and the classical theorem of Pringsheim.
It proves finite complex radius without identifying a discriminant zero with the pressure branch.
Consequently the first negative coefficient after the cancellation degree exists for every $m$; together with the preceding positive-window theorem it occurs at a finite even degree at least $6m$.
Exact computations give its location for $2\le m\le20$.
No sign density or large-$m$ scaling law is asserted.
\end{abstract}

\section{The sign theorem}
Use the cosine phase convention of the repository manuscript \cite{Repo}:
\[
 g_c(x)=\cos^2\!\pi(x-c),\qquad T(x)=2x\pmod1,\qquad
 p_m(c)=P(m\log g_c).
\]
All logarithms are natural. Set
\[
 P_m(t)=p_m(t/\pi)=\sum_{k\ge1}c_{m,k}t^{2k}.
\]
The parameter $m$ is a fixed positive integer, while $t$ is the phase variable scaled by $\pi$.
The source identifies pressure with the logarithm of the positive dominant eigenvalue of a finite Fourier matrix and proves real analyticity at every real phase, including the atomic phase $c=0$.
We reconstruct the spectral facts needed here in Section~\ref{sec:real}; their attribution remains with \cite{Repo}.

\begin{theorem}\label{thm:sign}
For every integer $m\ge2$, the Taylor radius $R_m$ of $P_m$ satisfies
$0<R_m\le4$. There are infinitely many indices $k$ with $c_{m,k}>0$ and infinitely many with $c_{m,k}<0$. Moreover,
\begin{equation}\label{eq:scale}
 \limsup_{k\to\infty}\bigl(\max\{c_{m,k},0\}\bigr)^{1/k}
 =
 \limsup_{k\to\infty}\bigl(\max\{-c_{m,k},0\}\bigr)^{1/k}
 =R_m^{-2}.
\end{equation}
\end{theorem}
Thus neither sign eventually disappears, and neither sign is exponentially smaller along all of its indices.
After zero coefficients are omitted, infinitely many sign changes occur.
This does not imply consecutive alternation or a density of either sign.
For $m=1$, pressure is identically zero and the conclusion is false.

The analytic mechanism is classical: a finite-radius nonnegative series has a singularity at its positive radius, by Pringsheim \cite{FS}. The application to this pressure family, rather than that general principle, is the result of this note.
The argument is independent of finite coefficient sampling.

\section{The finite matrix and the real axis}\label{sec:real}
Put $I_m=\{1-m,\ldots,m-1\}$ and
\[
 a_j=4^{-m}\binom{2m}{m+j},
\]
with $a_j=0$ outside $[-m,m]$. On the trigonometric polynomials
$\mathcal E_m=\operatorname{span}_{\C}\{e^{2\pi i r x}:r\in I_m\}$, the sum transfer operator
\[
 (\Lop_{m,c}f)(x)=\sum_{Ty=x}g_c(y)^m f(y)
\]
has matrix
\begin{equation}\label{eq:matrix}
 A_m(t)_{k,r}=2a_{2k-r}e^{-2i(2k-r)t},\qquad t=\pi c.
\end{equation}
Indeed, multiplication by $g_c^m$ gives frequencies of absolute value at most $2m-1$; summing over the two inverse branches removes odd frequencies and halves even frequencies. This proves invariance and \eqref{eq:matrix}.

At $t=0$, every column sums to one. The admissible values of $j=2k-r$ form one parity class, whose binomial mass is $1/2$, and all resulting indices $k$ belong to $I_m$.
Every nonzero entry has $|2k-r|\le m$. Consequently, for complex $t$,
\begin{equation}\label{eq:growth}
 \|A_m(t)\|_1\le e^{2m|\operatorname{Im}t|}.
\end{equation}

\subsection{The zero-weight Perron argument}
The following proof is the finite-dimensional argument of \cite[Sections 5--7]{Repo}; it is included to make the treatment of the weight's zero explicit.
For real $c$, $\Lop_{m,c}$ preserves the closed, pointed, solid, generating cone of nonnegative real trigonometric polynomials.
The finite-dimensional cone Perron principle supplies a nonzero $h\ge0$ with eigenvalue equal to the spectral radius $\rho$.
This principle follows, for example, by adding a small strongly positive rank-one map, applying the fixed-point theorem on a normalized cone base, and taking a limit.
For the inverse branches the two values of $g_c$ are $a$ and $1-a$, so
\[
 2^{1-m}\le \Lop_{m,c}1=a^m+(1-a)^m\le1.
\]
The supremum norm upper bound and iteration of the lower bound give
$2^{1-m}\le\rho\le1$.

\begin{lemma}\label{lem:tree}
Let $Z\subset\T$ be nonempty and finite. If
$T^{-1}Z\setminus\{s\}\subseteq Z$, then $Z=\{0\}$ and $s=1/2$.
\end{lemma}
\begin{proof}
Counting gives $2|Z|-1\le|Z|$, hence $|Z|=1$.
The remaining inverse image must be the same point, hence a fixed point of doubling.
The only such point is zero, whose other preimage is $1/2$.
\end{proof}

The weight vanishes only at $s=c+1/2$.
If $c\ne0\pmod1$, a nonempty zero set of $h$ would be finite and satisfy Lemma~\ref{lem:tree}, a contradiction. Thus $h>0$.
For an eigenfunction $\Lop f=\lambda f$ with $|\lambda|=\rho$, set
\[
 R=\max_x\frac{|f(x)|}{h(x)},\qquad H=R^2h^2-|f|^2.
\]
If $f\ne0$, then $R>0$, and $H$ is a nonnegative trigonometric polynomial with a nonempty zero set.
At a zero, equality holds in
\[
 |\lambda f(x)|
 \le\sum_{Ty=x}g_c(y)^m|f(y)|
 \le R\sum_{Ty=x}g_c(y)^m h(y)=R\rho h(x).
\]
Its zero set therefore propagates to both preimages except $s$.
Lemma~\ref{lem:tree} forces $H\equiv0$.
Writing $F=f/h$, equality in the triangle inequality now gives
\[
 F(y)=(\lambda/\rho)F(Ty)\quad(y\ne s),
\]
and continuity extends this to $s$.
At zero, $F(0)\ne0$ gives $\lambda=\rho$.
The dyadic preimages of zero are dense, so $F=F\circ T$ is constant.
Finally $\rho^{-1}\Lop$ contracts the norm $\max|f|/h$, excluding a nontrivial Jordan block.
Thus $\rho$ is algebraically simple and is the unique peripheral eigenvalue.

\subsection{The exceptional phase and analytic continuation}
At $c=0$, let $J$ have law $\Pr(J=j)=a_j$.
The transpose matrix acts on polynomial evaluations on $I_m$ by
\[
 (A_m(0)^{\mathsf T}Q)(r)
 =\E\!\left[Q\!\left(\frac{r+J}{2}\right)\,\middle|\,J\equiv r\pmod2\right].
\]
For every polynomial $U$ of degree less than $2m$,
$\sum_j(-1)^j a_j U(j)=0$, by the vanishing $2m$th finite difference.
The relevant conditional binomial moments therefore equal the unconditional moments.
On the evaluation basis $1,r,\ldots,r^{2m-2}$, the transpose is triangular with diagonal
\[
 1,2^{-1},\ldots,2^{-(2m-2)}.
\]
These are distinct, so the dominant eigenvalue one is simple at the exceptional phase too.

The implicit-function theorem for the characteristic polynomial now gives a positive real-analytic dominant eigenvalue at every real phase.
The local branches agree on overlaps by uniqueness.
Its logarithm is therefore real-analytic on the whole real line.
Reflection and periodicity give
$P_m(-t)=P_m(t)$ and $P_m(t+\pi)=P_m(t)$.

For completeness, evaluation at zero is a left eigenfunctional at the atomic eigenvalue one. It pairs nontrivially with its right eigenvector: otherwise that eigenvector would lie in the range of $A_m(0)-I$, creating a Jordan chain.
Thus the local analytic eigenfunction can be normalized by $h_{m,c}(0)=1$.
Evaluating its transfer equation at zero yields the source's delayed-feedback identity
\[
 \rho_m(c)=\cos^{2m}(\pi c)+\sin^{2m}(\pi c)\,h_{m,c}(1/2).
\]
For $m\ge2$ this proves
\begin{equation}\label{eq:nonconstant}
 P_m(t)=-m t^2+O(t^4).
\end{equation}
In particular the pressure germ is nonconstant.

\section{A finite complex radius and equal sign scales}
We first give a direct quantitative radius argument.
For every circle $|t|=r$ strictly inside the Taylor disk, the characteristic identity
\[
 \det\bigl(e^{P_m(t)}I-A_m(t)\bigr)=0
\]
holds by analytic continuation from zero. Equation~\eqref{eq:growth} gives
$\operatorname{Re}P_m(t)\le2mr$ on that circle.
Write $u(\theta)=\operatorname{Re}P_m(re^{i\theta})$ and $M=2mr$.
The mean of $u$ is $P_m(0)=0$. By the second Fourier coefficient and \eqref{eq:nonconstant},
\[
 -mr^2=2\frac1{2\pi}\int_0^{2\pi}u(\theta)e^{-2i\theta}\,d\theta
 =-2\frac1{2\pi}\int_0^{2\pi}(M-u(\theta))e^{-2i\theta}\,d\theta .
\]
As $M-u\ge0$, it follows that $mr^2\le2M=4mr$, hence $r\le4$.
Taking the supremum proves $0<R_m\le4$.

Equivalently, an entire pressure would satisfy
$\operatorname{Re}P_m(t)\le2m|t|$.
For an entire function $F(t)=\sum a_jt^j$ with $\operatorname{Re}F(t)\le C+L|t|$, the same harmonic Fourier estimate gives
\[
 |a_j|r^j\le2(C+Lr-\operatorname{Re}F(0)).
\]
Letting $r\to\infty$ makes $F$ affine.
An even such function is constant, contradicting \eqref{eq:nonconstant}.
Neither argument assumes that a chosen discriminant zero belongs to the pressure branch.

Now put
\[
 Q_m(z)=\sum_{k\ge1}c_{m,k}z^k,\qquad r_m=R_m^2.
\]
This series has finite positive radius $r_m$.
For each $x>0$, the positive real square root and the real-axis Perron logarithm give a local analytic function $P_m(\sqrt z)$.
These functions agree as they are continued along the positive axis.
At $x=r_m$, its local version agrees with the Taylor germ on a real interval $(r_m-\varepsilon,r_m)$ and hence, by the identity theorem, on their complex overlap.
Thus $Q_m$ continues analytically through its positive radius.

If the coefficients were eventually nonnegative, subtraction of a finite polynomial would give a nonnegative-coefficient series with radius $r_m$ and no singularity at $r_m$, contrary to Pringsheim's theorem \cite[Theorem IV.3(ii)]{FS}.
Applying the same argument to $-Q_m$ excludes eventual nonpositivity.
This proves both infinite-sign conclusions.

For the stronger statement define
\[
 Q_+(z)=\sum_{k\ge1}\max\{c_{m,k},0\}z^k,\qquad
 Q_-(z)=\sum_{k\ge1}\max\{-c_{m,k},0\}z^k.
\]
Their radii $r_+,r_-$ are at least $r_m$, and
$\min(r_+,r_-)=r_m$ by Cauchy--Hadamard.
If $r_+=r_m<r_-$, then $Q_+=Q_m+Q_-$ would continue through $r_m$, contradicting Pringsheim for $Q_+$.
The reverse inequality is equally impossible.
Therefore $r_+=r_-=r_m$, proving \eqref{eq:scale}.

\section{Both signs in each even degree class modulo four}
For real $y$, all nonzero entries of $A_m(iy)$ are strictly positive.
Its diagonal is positive, and there are positive entries between each pair of adjacent Fourier indices in both directions.
Indeed, for output $k=r\pm1$, the binomial index corresponds to $2k-r=r\pm2$, which remains in $[-m,m]$ whenever both adjacent indices belong to $I_m$.
The matrix is therefore irreducible and aperiodic, hence primitive.
Its Perron eigenvalue is positive and algebraically simple for every $y$.
Starting from the eigenvalue one at $y=0$, this supplies the analytic continuation of $e^{P_m(iy)}$ along the imaginary axis.
Its real logarithm gives a continuation of $P_m(iy)$ there.
It follows that $Q_m(z)$ is analytic at every negative real $z$ as well as every positive real $z$.
In particular, every nearest singularity of $Q_m$ is nonreal.

Define the two filtered series
\[
 E_m(w)=\sum_{j\ge1}c_{m,2j}w^j,\qquad
 O_m(w)=\sum_{j\ge0}c_{m,2j+1}w^j.
\]
They are the coefficients in pressure degrees $4j$ and $4j+2$, respectively.
\begin{theorem}\label{thm:filters}
Both $E_m$ and $O_m$ have finite positive Taylor radii.
Each has infinitely many strictly positive and infinitely many strictly negative coefficients.
Within each filtered series, its positive and negative parts have the same radius as that filtered series.
\end{theorem}
\begin{proof}
Let $D(t)=\operatorname{diag}_{k\in I_m}(e^{2ikt})$.
The matrix factorization and centered index set give
\[
 A_m(t)=D(-2t)A_m(0)D(t),\qquad
 \det A_m(t)=\det A_m(0)=2^{-(2m-2)(2m-1)/2}\ne0.
\]
Thus $A_m(t)^{-1}$ is entire and of exponential type: its adjugate entries are finite sums of products of exponential entries, and its denominator is a nonzero constant.

As germs at zero,
\[
 S_m(t):=P_m(t)+P_m(it)=2E_m(t^4),\qquad
 B_m(t):=P_m(t)-P_m(it)=2t^2O_m(t^4).
\]
The exponential of $S_m$ is an eigenvalue of $A_m(t)\otimes A_m(it)$.
The exponential of $B_m$ is an eigenvalue of $A_m(t)\otimes A_m(it)^{-1}$.
Both tensor matrices are entire and have norm at most $C_m e^{L_m|t|}$ for suitable fixed constants.
If $E_m$ were entire, then $S_m$ would be entire and its characteristic identity would extend globally.
Hence $\operatorname{Re}S_m(t)\le\log C_m+L_m|t|$, forcing $S_m$ affine by the harmonic estimate above, and constant by evenness.
The same argument applies to $B_m$ if $O_m$ were entire.

Neither germ is constant.
The difference has $[t^2]B_m=-2m$.
For $m\ge3$, delayed feedback begins in degree at least six, so
$[t^4]S_m=-m/3$ from $2m\log\cos t$.
For $m=2$, the cubic calculation in Section~\ref{sec:finite} gives
$[t^8]S_2=2\cdot6836/189\ne0$.
Thus both filtered radii are finite and positive.

For positive $w$, choose the positive fourth root $t=w^{1/4}$.
The real and imaginary Perron continuations just proved make both
$S_m(t)/2$ and $B_m(t)/(2t^2)$ analytic locally in $w$.
They glue to the corresponding filtered germs along the positive axis.
Pringsheim and the positive/negative decomposition used for Theorem~\ref{thm:sign} therefore apply separately to each filter.
\end{proof}
The two filtered radii need not be identified with each other by this argument.
If they are $r_E,r_O$, the precise relation to the original radius is
$\min(r_E,r_O)=R_m^4$.
No assertion about the frequency or spacing of these sign changes is required.

\section{First negative degrees: existence and exact finite values}\label{sec:finite}
The source proves $c_{m,m}=0$.
The preceding positive-window theorem \cite{Window} proves
$c_{m,k}>0$ whenever $m<k<3m$.
Theorem~\ref{thm:sign} therefore has the following consequence.
\begin{corollary}
For every integer $m\ge2$, the first post-cancellation negative degree
\[
 N_m=\min\{2k:k>m,\ c_{m,k}<0\}
\]
exists, is finite, and satisfies $N_m\ge6m$.
\end{corollary}

The source's rational Fourier recursion, with the full formal logarithm retained, gives the following exact finite values.
\begin{center}
\begin{tabular}{c|rrrrrrrrrr}
$m$&2&3&4&5&6&7&8&9&10&11\\
\hline
$N_m$&12&18&24&32&38&44&52&58&64&72
\end{tabular}

\smallskip
\begin{tabular}{c|rrrrrrrrr}
$m$&12&13&14&15&16&17&18&19&20\\
\hline
$N_m$&78&84&92&98&104&110&118&124&130
\end{tabular}
\end{center}
Every intervening coefficient is computed exactly, so each entry is a first negative degree, rather than only an upper bound.
The table does not establish a limit for $N_m/m$.

Here is the finite algorithm's normalization.
With $a=\tan t$, the row-polynomial gauge of the normalized phase matrix is
\[
 \operatorname{diag}_k\!\bigl((1-ia)^{m+k}(1+ia)^{m-k}\bigr)A_m(0).
\]
If its eigenvalue is $\lambda(a)$ near one, then
\[
 P_m(t)=2m\log\cos t+\log\lambda(\tan t).
\]
The checker solves its normalized eigenvector recursion over the rationals after the substitution $b=ia$ and verifies the original Fourier equation at every order.
Writing $\lambda(a)=1+\sum_{j\ge1}v_j a^{2j}$, it computes the full logarithm
$\log\lambda(a)=\sum_{j\ge1}q_j a^{2j}$ from
\begin{equation}\label{eq:log}
 q_j=v_j-\sum_{i=1}^{j-1}\frac{i}{j}q_i v_{j-i}.
\end{equation}
All feedback orders are included.
The tangent substitution is computed from $\tan'=1+\tan^2$.
Each case is checked through degree $10m$, beyond its displayed first negative degree.

As an independent comparison, a second program constructs the characteristic polynomial of the original phase matrix for $2\le m\le6$, solves its simple root at one directly as a series in $t^2$, and takes its logarithm.
All $105$ coefficients through the same respective degrees agree exactly.
For $m=2$, the characteristic equation is
\[
 \mu^3-(\cos2t+3/4)\mu^2+(5\cos2t/8+1/4)\mu-1/8=0,
 \qquad \mu(0)=1,
\]
and the first negative coefficient after cancellation is
\[
 [t^{12}]P_2(t)=-\frac{35360872}{93555}.
\]
The finite arithmetic supports the table; the infinite-sign theorem does not depend on it.

\section{Attribution, scope and reproducibility}
The real-phase Perron theorem, the finite Fourier matrix, the atomic spectrum and the cancellation are inputs from the repository manuscript \cite{Repo}.
The treatment of the weight's zero in Section~\ref{sec:real} reproduces its finite-dimensional argument.
Gohlke, Kesseb\"ohmer and Schindler \cite[Theorem 2.7]{GKS} establish the order-two real-phase regularity in the primary literature.
The complex harmonic estimate and Pringsheim's theorem are classical tools.
The result here is their application to all integer pressures, together with the equal exponential scales of the two signs.
The inspected source does not state this sign conclusion.
A targeted literature search found no earlier statement, but this is not an exhaustive priority claim.

The proof does not locate $R_m$, select a dominant algebraic branch point, give a density of signs, or determine the asymptotic location of $N_m$ as $m$ grows.
A real discriminant zero can involve two other eigenvalues and need not limit the pressure Taylor series.
The first-negative table should not be used as a substitute for that branch analysis.

The source bundle contains this editable manuscript, a build script, the exact rational checker and data for $m=2,\ldots,20$, the independent characteristic-polynomial checker, validation records and source hashes.
The proof has been checked by direct derivation and exact arithmetic; it is not externally refereed or formally verified.

\begin{thebibliography}{9}
\bibitem{Repo}
\emph{Integer Pressure and a Missing Taylor Coefficient: Analytic phase dependence and the sixth-moment optimum for generalized Thue--Morse products}.
ProveIt research manuscript, 28 September 2026; inspected 1 October 2026.
\href{https://github.com/VladimirReshetnikov/ProveIt/blob/63a7a325109ba611a1816b61dfd0eb072b896a7a/Analysis/FabiusFunction/docs/semi-formalized-research-frontiers/drafts/thue-morse/Thue_Morse_Integer_Pressure/article.tex}{Pinned repository source}.
Git blob \texttt{1444c01b4b30020f727172f0ee0060cab644f444}.

\bibitem{FS}
P.~Flajolet and R.~Sedgewick,
\emph{Analytic Combinatorics}, Chapter IV lecture-note draft,
Theorem IV.3(ii), printed page 15.
\url{https://algo.inria.fr/flajolet/Teach/fascII.pdf}.
See also the book, Cambridge University Press, 2009.

\bibitem{GKS}
P.~Gohlke, M.~Kesseb\"ohmer and T.~I.~Schindler,
\emph{Generalized Thue--Morse measures: spectral and fractal analysis},
arXiv:2509.22109v1, 2025.
\url{https://arxiv.org/abs/2509.22109}.

\bibitem{Window}
\emph{The Full Pressure Positivity Range at Every Integer Order}.
Companion research report prepared for Vladimir Reshetnikov with OpenAI,
1 October 2026; delivered as report 63.
It proves positivity at every even degree strictly between $2m$ and $6m$ for every integer $m\ge2$.
\end{thebibliography}
\end{document}
